% !TeX root = ../../main.tex
\section{从计算需求到执行成本的分层分析}
\subsection{算子、阶段与应用三个计时边界}
同一项优化在不同计时边界下可能产生不同结论。算子边界只覆盖 SpMV 或 Top-K 选择；阶段边界包含输入准备、算子调用和结果整理；应用边界还可能包含数据编码、索引构建与下游消费。本文主要讨论查询与求解的在线完整流程，并单独列出离线成本。明确边界的意义在于，只有实际落在关键路径中的成本变化才会影响用户可见的完成时间。

以两阶段检索为例，若仅测量 INT6 扫描内核，所观测到的是候选分数生成速率；若测量一次查询，还需要包含局部候选维护、合并与 FP32 重算。低比特扫描越快，其他阶段在总时间中的相对比例越高。因此，不能用低比特算子的加速比直接代表查询加速比。若不同阶段对连续查询形成流水，系统吞吐又由最慢阶段决定，单条查询延迟则仍包含其经过各阶段的时间。

对 CG，单轮 SpMV 用时下降并不充分。设参考实现执行 $N_0$ 轮，每轮平均时间 $t_0$；近似实现执行 $N_1$ 轮，每轮核心计算平均时间 $t_1$，每轮控制平均时间为 $t_c$，其他一次性在线成本为 $T_s$。则简化的求解加速比为
\begin{equation}
 S=\frac{N_0t_0}{N_1(t_1+t_c)+T_s}.
\end{equation}
该式没有假定每轮时间完全一致，而是以完整轨迹的平均量表示总和。它揭示两种约束：一是额外迭代抵消单轮节省，二是短任务中一次性成本的比例较大。真正的性能目标因而应是满足相同质量条件的总求解时间。

\begin{table}[htbp]\centering\caption{不同计时边界能够支持的结论}
\begin{tabularx}{\textwidth}{p{2.4cm}XX}\toprule
边界&可回答的问题&不能直接推出的结论\\\midrule
SpMV 内核&单位稀疏计算的执行效率&完整查询或求解的加速比\\
候选生成阶段&粗筛成本及候选覆盖情况&精排后的全部任务质量\\
在线查询&既定索引上的查询代价&包含索引更新的长期服务成本\\
完整迭代求解&达到停止条件的时间&包含策略训练的首次求解收益\\
有效使用周期&准备、更新与执行的总成本&不同更新模式下的同等收益\\\bottomrule
\end{tabularx}\end{table}

\subsection{数值位宽与实际传输量的关系}
考虑一个用于说明字段成本的理想记录：数值占 $b_v$ 位，列索引占10 位，暂不计行边界及对齐。将数值从32 位降到6 位，记录从42 位变为16 位，其字节数比例为 $16/42$，与 $6/32$ 不同。若再把512 位传输单元视为不可分割的包，容量分别为12 项与32 项，比例又受到取整影响。这些数值是字段预算示例，不替代第3 章的实际 Ultra-CSR 协议。

实际传输还存在固定开销。设有 $z$ 个保留非零元素，元数据总字节数为 $M$，每个元素平均物理字节数为 $b$，则
\begin{equation}
 B=M+zb.
\end{equation}
量化改变 $b$，剪枝改变 $z$，结构压缩可能改变 $M$。当 $M$ 占比很高时，继续减少数值位宽的边际收益会降低；当矩阵由大量极短行组成时，行指针、边界标记和包尾空闲尤其需要关注。仅报告数值数组大小会低估这些影响。

还应区别总传输量和传输速率。Random-CSR 只访问少量候选，但各行起点不连续，单次读取长度可能较短；Ultra-CSR 访问更多记录，却能形成长连续流。因此，两种路径的有效带宽并不必然相同。合理的两阶段设计允许高带宽流承担大量低成本筛选，让低带宽随机流只服务少量必要重算，而不是假定所有数据读取具有相同成本。

\subsection{工作量减少与任务碎片化的关系}
选择性计算把原本连续扫描拆为若干被选区域。假设一次查询访问 $s$ 个簇，第 $j$ 个簇有效数据为 $B_j$，每次启动的固定成本为 $\tau_j$，其有效传输速率为 $v_j$，串行路径上的读取成本可写为
\begin{equation}
 T_{\mathrm{read}}\approx\sum_{j=1}^{s}\left(\tau_j+\frac{B_j}{v_j}\right).
\end{equation}
该式不代表 HBM 控制器的精确排队模型，但说明相同总字节数可以对应不同延迟。把一个区域拆得更细，能够提高选择的针对性，也可能增加启动次数和地址处理。若这些事务可以重叠，固定成本的一部分被隐藏；是否隐藏则取决于并发请求能力、缓存和依赖关系。

因此，簇数增加具有双重作用：一方面改善候选区域的分辨率，另一方面增加代表评分、目录与短事务开销。聚类目标若只优化簇内相似度，就未必得到最快的物理执行。第4 章通过簇内连续布局和多核探测将这两种要求联系起来，但它仍需要根据具体数据分布选择工作点。

GPU 混合精度也存在类似现象。过细的动作粒度可以减少高精度元素，却可能产生更多小任务和更多调度信息。统一内核的作用是将异构数值选择限制在共享的结构遍历中，减少把动作差异放大为内核碎片的机会。由此可见，近似决策的粒度必须与硬件任务粒度共同设计。

\subsection{质量损失来源的逐层定位}
检索系统中的低召回可能来自量化、元素剪枝、簇选择、局部 Top-K 容量以及并列分数规则。若只报告最终 Recall，无法判断应该调整哪一层预算。可采用逐层恢复的分析方式：保留其余执行不变，仅把某层替换为精确或全量版本，观察结果变化。这里的“精确”必须指定参考数值类型与得分函数，而不是把任何浮点计算称为数学精确。

\begin{table}[htbp]\centering\caption{检索质量损失的定位方法与局限}
\begin{tabularx}{\textwidth}{p{2.2cm}XX}\toprule
可能来源&可采用的对照&对照仍不能消除的因素\\\midrule
数值量化&保留候选范围，改用 FP32 评分&范围遗漏和局部容量截断\\
非零剪枝&保留格式，恢复完整非零支撑&量化和选择机制的影响\\
簇级选择&保持评分方式，探测全部簇&数值扰动与局部截断\\
局部候选上限&汇集全体近似分数做全局选择&近似分数本身的排序扰动\\
结果并列规则&统一稳定行号作为次级键&其他误差来源\\\bottomrule
\end{tabularx}\end{table}

这些对照之间可能存在交互，因而各自改善的 Recall 不能简单相加。某一候选同时被簇选择和量化排除，单独修复其中一层时仍无法恢复它。完整分析需要观察候选集合的交、并与遗漏原因，而不只是报告几个百分比差值。

CG 的定位方式不同。需要在相同内核和相同运行时下比较精度策略，并保存残差轨迹、精度轨迹与每轮时间。对某轮发生的残差变化，可回放相同搜索方向分别执行 FP64 与混合精度 SpMV，分析局部扰动；但离线回放不包含扰动引起的后续状态改变，不能替代完整求解。两种分析分别回答单步误差和闭环后果，作用互补。

\subsection{从共性认识到章节技术选择}
上述分析表明，三项研究可以共享问题分解方式，但不能共享未经证明的性能或误差结论。第3 章选择较粗的阶段分工，因为候选得分可以在筛选后重新计算；第4 章选择簇级范围控制，因为数据区域具有可利用的相关性；第5 章选择反馈精度控制，因为当前误差会改变后续求解状态。

三者也具有不同的撤销能力。候选进入重算集合后，可以恢复其高精度分数，但被提前删除的候选无法由精排自动恢复。CG 中已经发生的低精度更新通常不能靠下一轮简单升精度完全撤销，可能需要残差校正或重启。由此决定了质量检查位置和保守预算不同，不能将检索中的高精度重算直接类比为迭代中的无代价纠错。

从执行角度看，FPGA 能够围绕固定数据协议配置算术位宽与排序逻辑，适合将阶段结构固化为流水；GPU 则依靠共享内核和批量执行吸收动态动作，适合在通用并行资源上控制运行时开销。本文依据任务语义选择硬件执行方式，使决策粒度、存储布局和并行组织相互匹配。
